\section{Conclusion}
\label{sec:conclusion}

A constraint inferred from data describes what the data happens to satisfy, whereas a constraint contained in governing text describes what the data is required to satisfy. This distinction  has received limited attention in database research, even though many published datasets are accompanied by documents that define data quality criteria at some level. 

In this paper, we show that the gap between data and normative text can be addressed automatically. Our results indicate that lightweight classification is sufficient to detect  most checkable requirements, and that a structured multi-step pipeline grounded in both text and schema can translate them into executable constraints while avoiding hallucinated predicates and literals. 

This enables a new form of validation: \system\ can reveal compliance failures even when they do not manifest as discoverable pattern in the data. In such settings, the normative document becomes  part of the data specification, and keeping that specification aligned with evolving text emerges a concrete challenge for data management.
\altic{[R2-O6.1] Reconhecer como limitação que os datasets são pequenos em nº de sentenças e que só SISU se aproxima do cenário real. Adicionar discussão de escalabilidade para documentos legais típicos (centenas de sentenças, cláusulas aninhadas): custo linear por sentença já é argumento a favor; discutir segmentação de cláusulas aninhadas e o modo document-level como rota, e apontar isso como trabalho futuro.}
\altic{[R1-O3] Responder à reusabilidade sob dados dinâmicos: como a constraint é extraída do TEXTO (não dos dados), ela permanece válida enquanto a norma não muda, e pode ser re-executada sobre novas instâncias/atualizações sem re-descoberta. Conectar com a observação final sobre manter a especificação alinhada ao texto em evolução. Distinguir ``dados mudam'' (constraint estável) de ``norma muda'' (re-extração).}
